\mode*
\section{Conclusions}
\label{conclusions}

We have argued, from variation theory, that debugging is self-directed
learning: a bug marks a critical aspect the programmer has not yet
discerned, and it is the bug itself that opens the dimension---by
contrasting the actual behaviour with what was expected---while fixing it
requires generating the remaining patterns of variation, contrast,
generalisation and fusion, for oneself.
Since Marton connects the deep approach to learning to exactly this
self-generated variation, we hypothesise that debugging skill and the deep
approach are the same underlying ability.
Let us conclude by restating each research question and answering it as
far as our findings allow.

\rqpatterns*
For experienced programmers we can already answer this: in the pilot
(\cref{pilot}), all three patterns were observable in the debugging
episodes, and the debuggers could not explain the phenomenon until they
had generated all of them.
The worked example (\cref{debugging}) shows what the patterns look like in
an episode.
For students the answer awaits the prgi26 data; whichever way it falls,
the coded episodes will give a catalogue of which patterns novices
generate spontaneously, in which order, and where they stop
(\cref{results}).

\rqdeep*
No empirical answer yet: this is what \cref{h:deep} is for, and the answer
comes with the autumn 2026 data.
What we can conclude already is theoretical: by \citeauthor{NCOL}'s
reinterpretation, the deep approach \emph{is} the disposition to create
the necessary patterns of variation (\cref{approaches}), so the question
is whether that disposition shows in the recorded debugging.
A negative finding would therefore strike at the reinterpretation---or at
our operationalisation---not merely at this study's design.

\rqsuccess*
Also pending \cref{h:success} and the data.
The pilot hints at the association: the explanation did not come until the
patterns were complete.
But the pilot involved no approach measure and no failed episodes, so it
is a hint, no more.

\rqexperts*
The pilot gives a preliminary answer for the expert side: expert debuggers
exhibit all three patterns, varying example after example until the
explanation arrives.
The students' side of the comparison awaits the study; the differences
between the two will point out precisely what debugging instruction needs
to provide (\cref{analysis}).

\rqgenerative*
Pending the within-cohort comparison in prgi26 (\cref{h:generative}).
Both outcomes are informative: if the material users generate more
patterns, generative learning transfers to debugging; if the self-studiers
generate just as many while scoring deeper on the questionnaire, the
ability comes with the approach rather than with the instruction
(\cref{results}).

If the hypotheses hold, debugging instruction becomes instruction in
self-directed variation: teach the students what to vary and what to hold
invariant, and---by generative learning---teaching the course content
through patterns of variation should improve their debugging even without
explicit debugging instruction.
The study runs in the autumn of 2026.

\begin{frame}<presentation>
  \begin{idea}
    \begin{itemize}
      \item Debugging \(=\) learning without a teacher.
      \item Deep approach \(=\) generating one's own patterns of variation.
      \item Hypothesis: they are the same skill.
      \item prgi26 + learnlog + R-SPQ-2F will tell.
    \end{itemize}
  \end{idea}
\end{frame}
